% Compiles on its own (pulls ../main.tex's preamble via subfiles) or as part of the book.
\documentclass[../main.tex]{subfiles}
\begin{document}
\noindent\fbox{\parbox{0.94\textwidth}{\textbf{Historical, quarantined
experiment.} This chapter selected a recipe using the official test set and is
not part of the corrected study. See \texttt{../../../EXPLANATION.md}.}}\par\medskip

\chapter{Choosing the Learning Rate}

The machine can train (Ch.~7). Before we spend real GPU-hours on the experiments that produce the
study's numbers, we owe one measurement: the learning rate. This chapter is that measurement --- how
we ran it, what it returned, what the returned numbers mean when read against the published
literature, and one uncomfortable thing the exercise revealed about our recipe that we did not go
looking for.

\section{Why this is a measurement and not a lookup}

The learning rate is the size of the step the optimiser takes downhill after every batch (Ch.~7,
the foggy-hillside picture). It is the single most consequential number in training: too large and
the descent bounces across the valley and never settles or diverges outright; too small and it
crawls, wasting hours.

For most projects the learning rate is a \emph{lookup}: you copy a value from a model someone
already tuned. We cannot do that, and the reason is Soham's from-scratch directive. A pretrained
model arrives with weights already in a good region, so its safe learning rate is known and
transferable. Ours starts from noise. Nobody has tuned an 8.94M-parameter YOLOX trained from
scratch on 25{,}086 KAIST thermal frames at batch~8. So the learning rate is an \emph{empirical
unknown for our exact setup}, and the only honest way to get it is to try several and measure.
That is a \textbf{learning-rate sweep}.

\section{How the sweep was run, and why single-seed}

We trained one short run per candidate rate and evaluated each on the held-out test set, reporting
MR$^{-2}$ on All/Day/Night. The candidates were spaced roughly logarithmically --- $0.04, 0.02,
0.01, 0.005, 0.0025, 0.00125$ across two attempts --- because learning rate acts multiplicatively;
equal \emph{ratios}, not equal \emph{differences}, are what matter.

Each run used a \textbf{single seed}, on purpose. This is a subtle methodological point worth
stating plainly, because it looks like it contradicts Soham's ``5 seeds on every experiment''
directive and does not. The five-seed rule exists to put error bars on a \emph{result} --- to
separate a real improvement from initialisation luck. The sweep is not producing a result; it is
making a coarse \emph{choice} among learning rates that differ by factors of two. A choice that
coarse does not need error bars: if two rates are within seed-noise of each other, either is fine,
and if one is clearly better, one seed already shows it. Spending five seeds per candidate here
would multiply the cost fivefold to sharpen a decision that does not need sharpening. The seeds go
where they matter: the final baselines.

\section{The cliff: what divergence actually looks like}

Two of the candidates never produced a usable model, and how they failed is worth seeing rather
than being told. At $\text{lr}=0.02$ the objectness loss did not descend --- it \emph{ran away}:

\begin{center}
\includegraphics[width=0.82\textwidth]{lr_divergence.png}
\end{center}

Read the shape. For the first few hundred steps the loss is merely noisy (the spike near step 840
is a single bad batch, recovered immediately). Then, past step~1000, it stops recovering and climbs
--- 12, 90, 257, 402, 664, 989 --- each step's oversized update making the next gradient larger, a
runaway feedback loop. The step size was large enough that every correction overshot into worse
territory. This is divergence: not slow learning, but the loss actively exploding.

It also \emph{crashed the process}, and the mechanism is a lesson in itself. When the loss explodes
the model's predicted numbers become \texttt{NaN} (not-a-number). Inside SimOTA's cost computation
sits a cross-entropy call that carries a \emph{device-side assertion} its input lies in $[0,1]$; a
\texttt{NaN} trips it. Crucially, a device-side assert does not raise a normal Python error --- it
\textbf{corrupts the entire CUDA context for the process}, so every later operation fails too. In
our first sweep this took down not just $\text{lr}=0.02$ but the three healthy rates queued behind
it. Only the first failure was real; the rest was collateral from a poisoned GPU state.

\begin{lesson}
Divergence must be made survivable, or one bad hyperparameter costs a whole sweep. We fixed this in
three layers: (1) sanitise \texttt{NaN}s before the cross-entropy so divergence stays a normal
Python event instead of a device assert; (2) detect a runaway loss on a \emph{smoothed} average
(an EMA over a grace period, not a single batch --- a first attempt used a single-batch threshold
and false-fired at step~22, because the per-image loss legitimately spikes when a batch holds only
one assignable pedestrian); and (3) train each learning rate in its own \emph{subprocess}, so even
a hard GPU fault costs one run, not all of them. The dashed line in the figure is that EMA guard.
\end{lesson}

There is a second, quieter finding hiding in the cliff. $\text{lr}=0.04$ --- \emph{larger} than the
0.02 that diverged --- actually survived all eight epochs (though it never settled, spiking and
dropping positives to the end). A larger rate surviving while a smaller one blows up is not a
contradiction: this model sits close to its stability edge, and which side of the edge a run lands
on depends partly on the luck of its random initialisation. That fragility has a real consequence
we return to in the recipe decision.

\section{The result}

The four rates that produced usable models:

\begin{center}
\includegraphics[width=0.9\textwidth]{lr_sweep_mr.png}
\end{center}

\begin{center}
\begin{tabular}{lcccl}
\toprule
lr & MR$^{-2}$ All & Day & Night & \\
\midrule
\textbf{0.01} & \textbf{36.56} & 43.47 & 22.61 & winner \\
0.0025 & 38.67 & 44.71 & 26.11 & \\
0.005 & 38.98 & 44.81 & 26.89 & \\
0.00125 & 43.21 & 50.94 & 27.35 & worst \\
\bottomrule
\end{tabular}
\end{center}

The curve is clean and monotonic: performance improves as the rate rises toward $0.01$, then the
cliff above it. The winner, $0.01$, beats the next-best by $2.1$ MR points --- and that margin is
\emph{real}, not noise. The bottom three cluster within four points, and the $0.005$/$0.0025$ pair
differ by $0.3$, which is about the size of single-seed jitter. Against that noise scale, a
$2.1$-point separation is a genuine signal. The optimum sits at or just below $0.01$.

\section{Reading the numbers honestly}

Two things in the table are worth dwelling on.

\textbf{Night beats Day, by a lot} ($22.6$ vs $43.5$). This is the thermal sensor's physics falling
straight out of the model. A thermal camera sees heat: at night a warm body blazes against a cool
background; in daytime the sun warms pavement, walls and cars to near body temperature and the
person no longer stands out thermally. The whole KAIST literature reports thermal-only dominating
at night, and here our from-scratch model has rediscovered it unaided. Had Day beaten Night we would
suspect a swapped label or a broken split. It did not --- the model learned real physics, which is
the strongest sanity check we have that the pipeline is sound end to end.

\textbf{The absolute number is mediocre, and that is expected.} Placed against the published record
(from \texttt{docs/LITERATURE.md}, thermal-only on KAIST):

\begin{center}
\begin{tabular}{lcccl}
\toprule
method & All & Day & Night & note \\
\midrule
FasterRCNN-T (2016) & 47.6 & 50.1 & 40.9 & original anno, hand-tuned era \\
\textbf{ours} & \textbf{36.6} & 43.5 & 22.6 & from scratch, 8 epochs \\
modern thermal-only & $\sim$19.0 & $\sim$24.6 & $\sim$7.8 & pretrained backbone \\
\emph{full-fusion SOTA} & 7--8 & --- & --- & for scale only \\
\bottomrule
\end{tabular}
\end{center}

We sit between the crude 2016 detector and a modern pretrained one --- exactly where a small,
from-scratch, lightly-trained model belongs. Better than 2016; well short of modern. Nothing is
broken, nothing is magical.

\section{The uncomfortable discovery: we have not chosen the epoch budget}

Benchmarking forces an observation the sweep was not designed to produce. The gap between our
$36.6$ and modern $\sim$19 is mostly \emph{not architecture} --- it is \textbf{epochs and
pretraining}. YOLOX trains for on the order of 300 epochs from scratch on COCO. Our runs saw each
image \emph{eight} times. Eight epochs was a deliberate choice \emph{for the sweep} --- just enough
to rank learning rates --- but that means the number is a sweep artifact, not a recipe.

This matters far beyond a nicer headline, and it is the real payoff of thinking about the result
rather than just filing it:

\begin{lesson}
An undertrained baseline invalidates the whole study, silently. PRISM's central question is whether
fusion beats single-modality. If the thermal-only baseline is undertrained and the fusion model is
trained to convergence, then ``fusion wins'' is uninterpretable --- you cannot tell whether fusion
\emph{helped} or merely \emph{trained more effectively}. A fair comparison requires baseline and
fusion to reach comparable convergence. So the epoch budget is not cosmetic; it is a controlled
variable, and it must be fixed \emph{before} the reference baselines are run. We locked the learning
rate here; we have not yet locked epochs, and that is the next thing to settle.
\end{lesson}

\section{The decision, with margin}

We lock \textbf{lr $=0.01$}. We considered probing higher (0.012--0.015, since $0.01$ is the top of
the tested range), and rejected it. The reasoning is stability, not just the $2.1$-point win:
$0.01$ sits a comfortable $2\times$ below the $0.02$ divergence cliff, and because this model lives
near its stability edge, the five-seed protocol will run five different initialisations at this
rate. A rate pushed toward the cliff risks some of those seeds diverging and wrecking the noise
floor. A recipe chosen at the edge is a fragile recipe; a marginal fraction of an MR point is not
worth trading the robustness of every downstream run for. Locked at $0.01$, with margin.

\medskip\noindent\textbf{Locked so far:} learning rate $=0.01$. \textbf{Still open:} the epoch
budget (Ch.~9), and the annotation sets that make our numbers comparable to the leaderboard (the
on-disk ``sanitized'' set is a duplicate of the ``new'' set --- a real download is pending; see the
decisions log). Only once epochs are fixed and the correct annotations are installed do the
five-seed baselines get spent.

\end{document}
